% answers_short.tex —— 2014 年·全国大纲卷·理科（广西） · 简短答案速查
% 由 tools/make_short_answers.py 自 answers.tex 的 \gksolhead/\gkanswer 提取，
% 并与 2014/2014_数学真题_大纲理_答案要点.md 逐题比对后生成，勿手改（改源 MD 后重跑）。

\begin{gknotice}
\textbf{使用说明}：本节只给最终答案便于快速核对，完整推导见解析版。答案为编写者独立推导，\textbf{非官方标准答案}。
\end{gknotice}

\gkpart{选择题}
\begin{center}\begin{tabular}{ccccccccccccc}
\toprule
题号 & 1 & 2 & 3 & 4 & 5 & 6 & 7 & 8 & 9 & 10 & 11 & 12 \\
\midrule
答案 & D & B & C & B & C & A & C & A & A & C & B & D \\
\bottomrule
\end{tabular}\end{center}

\gkpart{填空题}
\begin{itemize}[leftmargin=2.6em]
  \item[\textbf{13.}] $70$。通项为 $\mathrm C_8^k(-1)^k x^{8-\frac{3k}{2}}y^{\frac{3k}{2}-4}$，令指数为 $2$ 得 $k=4$，系数 $\mathrm C_8^4=70$。
  \item[\textbf{14.}] $5$。可行域的三个顶点为 $(1,1)$、$(-1,-1)$、$\left(2,\dfrac12\right)$，对应 $z$ 值为 $5$、$-5$、$4$，最大值为 $5$。
  \item[\textbf{15.}] $\dfrac43$。设夹角为 $\theta$，则 $\sin\dfrac{\theta}{2}=\dfrac{\sqrt2}{\sqrt{10}}=\dfrac1{\sqrt5}$，$\tan\dfrac{\theta}{2}=\dfrac12$，故 $\tan\theta=\dfrac43$。
  \item[\textbf{16.}] $(-\infty,2]$。在 $\left(\dfrac\pi6,\dfrac\pi2\right)$ 上 $\cos x>0$，$f'(x)\leqslant0$ 等价于 $a\leqslant4\sin x$ 恒成立，右端下确界为 $4\sin\dfrac\pi6=2$。
\end{itemize}

\gkpart{解答题}
\begin{itemize}[leftmargin=2.6em]
  \item[\textbf{17.}] $B=\dfrac{3\pi}{4}$。由正弦定理得 $3\sin A\cos C=2\sin C\cos A$，即 $3\tan A=2\tan C$，故 $\tan C=\dfrac12$，$\tan(A+C)=1$，$A+C=\dfrac{\pi}{4}$。
  \item[\textbf{18.}] (1) $a_n=13-3n$。由 $S_n\leqslant S_4$ 得 $a_4\geqslant0$、$a_5\leqslant0$，即 $-\dfrac{10}{3}\leqslant d\leqslant-\dfrac52$，又 $a_2=10+d$ 为整数，故 $d=-3$；(2) $T_n=\dfrac{n}{10(10-3n)}$。由 $b_n=\dfrac13\left(\dfrac1{10-3n}-\dfrac1{13-3n}\right)$ 裂项相消得 $T_n=\dfrac13\left(\dfrac1{10-3n}-\dfrac1{10}\right)$
  \item[\textbf{19.}] (1) 见解析版；(2) $\arccos\dfrac14$。由线面距离 $\sqrt3$ 得 $h^2=3(2-t)^2$，结合 $h^2=4t-t^2$ 及 $D$ 在 $AC$ 上得 $t=1$、$h=\sqrt3$，两半平面法向量夹角的余弦为 $\dfrac14$
  \item[\textbf{20.}] (1) $0.31$。$P(X=4)=0.06$，$P(X=3)=0.25$，故 $P(X\geqslant3)=0.31$；(2) $2$。四人是否使用设备相互独立，期望等于各人概率之和 $0.6+0.5+0.5+0.4=2$
  \item[\textbf{21.}] (1) $y^2=4x$。由 $|PQ|=\dfrac8p$、$|QF|=\dfrac8p+\dfrac p2$ 及 $|QF|=\dfrac54|PQ|$ 解得 $p=2$；(2) $l:y=x-1$ 或 $l:y=-x+1$。设 $l:x=my+1$，由 $A$、$M$、$B$、$N$ 四点共圆得四纵坐标之和为 $0$，即 $4m-\dfrac4m=0$，故 $m=\pm1$
  \item[\textbf{22.}] (1) 当 $1<a<2$ 时在 $(-1,a(a-2))$、$(0,+\infty)$ 上递增，在 $(a(a-2),0)$ 上递减；当 $a=2$ 时在 $(-1,+\infty)$ 上递增；当 $a>2$ 时在 $(-1,0)$、$(a(a-2),+\infty)$ 上递增，在 $(0,a(a-2))$ 上递减；(2) 见解析版
\end{itemize}
